%% kapa: exam=1 round=35 year=2024 period=2 form=A subject=accounting subject_ko=회계학 seq=40 number=80
%% source: Q-Net 감정평가사 기출문제 (cst00309) artlSeq=5213229 file=raw/1008/5213229/2교시 문제지 원본.pdf
%% conversion: kordoc 4.24.51 -> tools/gen1.py (auto); review=proofread-sonnet; flags=none
\begin{question}{80}
최신의 관리회계기법에 관한 설명으로 옳지 않은 것은?
\begin{choices}
  \item 목표원가는 목표가격에서 목표이익을 차감하여 결정한다.
  \item 카이젠원가계산은 제조이전단계에서의 원가절감에 초점을 맞추고 있다.
  \item 균형성과표는 조직의 전략과 성과평가시스템의 연계를 강조하고 있다.
  \item 품질원가의 분류에서 내부실패원가는 불량품의 재작업원가나 폐기원가 등을 말한다.
  \item 제품수명주기원가계산은 단기적 의사결정보다는 장기적 의사결정에 더욱 유용하다.
\end{choices}
\end{question}
\answer{2}
